% !TEX root = ../main.tex

% PATTERN: \nameref{label} — cross-reference a una sezione per nome, non numero.
% Esempio: "come illustrato in \nameref{esperimenti}" produce "come illustrato in Esperimenti".

\chapter{Discussione e Sviluppi Futuri}\label{chap:discussion}

Il Capitolo~\ref{chap:experiments} ha risposto a una domanda circoscritta: APIGuard Assurance produce i risultati dichiarati su un target reale? I dati confermano. Rimane aperta una domanda di portata diversa: entro quali confini strutturali quella risposta vale, e come il sistema si inserisce in contesti d'uso più ampi di quello sperimentale?

Questo capitolo esamina quei confini in tre direzioni. La prima è critica: i limiti strutturali del paradigma contract-driven, ciascuno radicato in una proprietà architetturale specifica e non eliminabile con un refactoring locale, vengono discussi senza attenuarli (Sezione~\ref{sec:limiti}). La seconda è applicativa: le condizioni concrete in cui il tool è integrabile in pipeline \gls{CI/CD} industriali e i pattern di deployment che ne massimizzano il valore (Sezione~\ref{sec:devsecops}). La terza è prospettica: le traiettorie di sviluppo distinte in estensioni operative che seguono direttamente dall'architettura attuale, e direzioni di ricerca che pongono problemi aperti e richiedono validazione sperimentale indipendente (Sezione~\ref{sec:sviluppi-futuri}).


% ─────────────────────────────────────────────────────────────────────────────
\section{Analisi Critica e Limiti del Paradigma Contract-Driven}%
\label{sec:limiti}
% ─────────────────────────────────────────────────────────────────────────────

Ogni proprietà architetturale che garantisce un vantaggio specifico impone, per costruzione, una classe di situazioni in cui quel vantaggio diventa vincolo. I limiti discussi in questa sezione non sono anomalie correggibili per via implementativa: sono conseguenze dirette di scelte deliberate, e come tali richiedono di essere nominati con precisione piuttosto che attenuati.


\subsection{Il Paradosso del Ground Truth (D3.P2)}\label{subsec:limite-ground-truth}

Come stabilito nella Sezione~\ref{subsec:openapi-sot}, la specifica OpenAPI è l'unico oracolo formale della pipeline: delimita la superficie di attacco, guida la costruzione dei probe e fornisce il contratto rispetto al quale le risposte vengono interpretate. Il vantaggio è la precisione semantica: dove un fuzzer cieco produce \texttt{HTTP~400} a causa di payload strutturalmente invalidi, APIGuard Assurance costruisce richieste conformi al contratto e sposta l'osservazione dalla sintassi alla semantica del controllo di sicurezza. Come discusso nella Sezione~\ref{subsec:openapi-sot}, questo è precisamente il \emph{combinatorial wall} che il paradigma contract-driven abbatte.

Il prezzo è la dipendenza dalla qualità di ciò che si assume come oracolo. Una specifica che non riflette endpoint aggiunti dopo l'ultima pubblicazione restringe la superficie di attacco osservabile senza che il sistema possa rilevare l'incoerenza: l'assenza di un endpoint dall'assessment è indistinguibile dalla sua assenza dalla specifica. Una specifica costruita per escludere endpoint sensibili produce un assessment formalmente corretto su una superficie artificialmente ristretta. In entrambi i casi il tool non dispone di un meccanismo per distinguere l'assenza di vulnerabilità dall'assenza di copertura; lo specification drift tra documentazione e implementazione reale è una delle cause principali di copertura incompleta nel testing contract-driven, indipendentemente dallo strumento adottato. Qualsiasi approccio che utilizzi una specifica formale come oracolo eredita la stessa dipendenza: la correttezza dell'assessment è necessariamente condizionale alla correttezza del contratto. La traiettoria di ricerca che affronta strutturalmente questo paradosso, l'inferenza automatica della specifica dall'analisi del traffico reale tramite tecniche di \gls{ML}, è discussa nella Sezione~\ref{sec:sviluppi-futuri} come direzione aperta, non come soluzione disponibile.

In termini operativi, la mitigazione praticabile è procedurale: l'operatore che esegue l'assessment è responsabile di verificare che la specifica fornita sia aggiornata e completa prima di interpretare i risultati. Il tool segnala parzialmente questa dipendenza producendo un \texttt{SKIP} con motivazione esplicita quando uno scope di test è vuoto, il che indica all'operatore che quel dominio non ha endpoint applicabili nella specifica fornita; non può tuttavia distinguere se lo scope è vuoto perché la specifica è corretta e quegli endpoint non esistono, o perché la specifica è incompleta e quegli endpoint esistono ma non sono documentati.


\subsection{La Tensione dell'Astrazione (D1.P6)}\label{subsec:limite-astrazione}

La Sezione~\ref{subsec:gateway-adapter} ha introdotto \texttt{BaseGatewayAdapter} come interfaccia astratta verso il piano di configurazione del gateway: i test \texttt{WHITE\_BOX} interrogano il gateway attraverso questa interfaccia, e aggiungere il supporto per un nuovo prodotto richiede una nuova implementazione concreta senza toccare i test che la utilizzano. Il vantaggio è la portabilità; il vincolo conseguente è che l'interfaccia può esporre solo ciò che è concettualmente generalizzabile a tutti i gateway che la implementano.

Un esempio rende la tensione concreta. Kong supporta l'ordinamento esplicito dei plugin tramite il campo \texttt{ordering} della sua Admin \gls{API}: un plugin di autenticazione eseguito dopo un plugin di trasformazione del payload riceve dati già alterati, con potenziali implicazioni sulla sicurezza. Verificare questa condizione richiede metadati Kong-specifici privi di equivalente strutturale in altri gateway; un test che li utilizzasse dovrebbe dipendere da \texttt{KongGatewayAdapter} direttamente, rompendo l'agnosticismo. Il design attuale non percorre questa strada.

La tensione non è necessariamente irrisolvibile. Una direzione di sviluppo percorribile è l'estensione della gerarchia dell'adapter: \texttt{BaseGatewayAdapter} rimane l'interfaccia comune con i metodi generalizzabili, ma classi intermedie come \texttt{BaseKongAdapter} o \texttt{BaseAWSAdapter} potrebbero esporre metodi vendor-specifici per i test che li richiedono, senza impattare i test che operano sul livello comune. Questa struttura a tre livelli rifletterebbe la stessa logica della gerarchia dei connector descritta nella Sezione~\ref{subsec:connector-hierarchy}. Un deployment che usa esclusivamente Kong potrebbe adottare test scritti contro \texttt{BaseKongAdapter}; un deployment multi-vendor userebbe solo i test del livello comune. La ricerca necessaria riguarda la specifica del contratto di compatibilità tra i livelli, non la fattibilità tecnica dell'approccio.


\subsection{La Dipendenza dall'Admin API (D4.P5)}\label{subsec:limite-admin-api}

La degradazione controllata descritta nella Sezione~\ref{subsec:graceful-degradation} garantisce che l'assenza della Admin \gls{API} non blocchi la pipeline, ma non elimina le conseguenze sulla copertura. In ambienti cloud managed, come Amazon \gls{API} Gateway, Azure \gls{APIM} o Google Cloud Endpoints, il piano di configurazione non è esposto tramite un'Admin \gls{API} interrogabile direttamente: è accessibile solo attraverso le \gls{API} del cloud provider, con modelli di autenticazione specifici per piattaforma. Configurare \texttt{gateway\_adapter} per questi ambienti richiede un adapter dedicato per ciascun provider, attualmente non implementato.

La copertura si riduce ai domini \texttt{BLACK\_BOX} e \texttt{GREY\_BOX}: D0, D1, D2, D3, D4 parzialmente, D7. Il Dominio~6 (Configuration \& Hardening) e le verifiche \texttt{WHITE\_BOX} di D4 restano escluse. La riduzione non è silenziosa: ogni \texttt{SKIP} riporta nel report la motivazione esplicita, rendendo la copertura effettiva immediatamente leggibile.

La prassi industriale che affronta questo limite è l'accordo di accesso privilegiato: nell'ambito di un assessment formale, il team di sicurezza ottiene accesso read-only alle \gls{API} di gestione del cloud provider per la durata dell'analisi, tipicamente con credenziali temporanee e scope limitato alle risorse del progetto. Questo modello presuppone un accordo esplicito con il provider o con il team di governance interna, e garantisce che la lettura della configurazione avvenga senza possibilità di modifica. Con quell'accesso disponibile, \texttt{gateway\_adapter} nel \texttt{config.yaml} viene popolato con le credenziali e l'endpoint dell'adapter specifico per il provider, e la copertura si estende ai domini \texttt{WHITE\_BOX} senza modifiche al codice. La condizione è che quell'adapter esista: la sua implementazione per i cloud provider principali è una delle estensioni operative discusse nella Sezione~\ref{subsec:roadmap-operativa}.



















% ─────────────────────────────────────────────────────────────────────────────
\section{Applicabilità Industriale e Integrazione DevSecOps}%
\label{sec:devsecops}
% ─────────────────────────────────────────────────────────────────────────────

Le proprietà di riproducibilità e degradazione controllata discusse nei Capitoli~\ref{chap:methodology} e~\ref{chap:implementation} non sono solo garanzie interne al tool: definiscono le condizioni concrete in cui APIGuard Assurance è integrabile come componente di una pipeline \gls{CI/CD} industriale. Questa sezione descrive quelle condizioni: la semantica del canale di comunicazione con la pipeline, i due profili di esecuzione disponibili e la variabilità di copertura nei diversi contesti organizzativi.


\subsection{Semantic Exit Codes come Canale di Comunicazione con la Pipeline (D6.P1)}\label{subsec:exit-codes}

Il contratto tra APIGuard Assurance e la pipeline che lo ospita è espresso da quattro valori interi. Quando il processo termina, l'exit code riflette lo stato aggregato del \texttt{ResultSet} calcolato da \texttt{report/builder.py}: \texttt{0} indica che nessun test ha prodotto \texttt{FAIL} o \texttt{ERROR}; \texttt{1} indica che almeno un test ha prodotto \texttt{FAIL}; \texttt{2} indica che nessun test ha prodotto \texttt{FAIL} ma almeno uno ha prodotto \texttt{ERROR}; \texttt{10} indica un errore infrastrutturale bloccante nelle fasi di startup (Fase~1, 2 o~4), ovvero il processo non ha mai raggiunto l'esecuzione dei test. La priorità tra i codici è ordinata: \texttt{FAIL} prevale su \texttt{ERROR}, quindi un assessment con entrambe le condizioni restituisce~\texttt{1}.

La separazione tra \texttt{1} e \texttt{10} è il punto architetturalmente rilevante. Da una pipeline automatizzata, "almeno un test ha trovato una violazione" e "il tool non si è avviato" sono condizioni con diagnosi e azioni correttive completamente diverse; trattarle con lo stesso valore scarica sull'operatore l'onere di leggere i log per distinguerle. Con exit code semanticamente distinti, un alert su \texttt{10} indica un problema infrastrutturale nel deployment della pipeline, mentre un alert su \texttt{1} indica che l'assessment è stato completato e ha rilevato almeno una garanzia violata.

L'integrazione in sistemi di \gls{CI/CD} come GitHub Actions o GitLab \gls{CI}\footnote{GitHub Actions e GitLab CI sono piattaforme di integrazione continua che eseguono pipeline di build, test e deployment automatizzato al variare del codice sorgente. Entrambe trattano qualsiasi exit code diverso da zero come fallimento del passo, bloccando l'avanzamento della pipeline.} non richiede wrapper aggiuntivi: il check sull'exit code è il meccanismo nativo. Un passo che esegue \texttt{apiguard run -{}-config config.yaml} blocca la pipeline se il tool restituisce exit~\texttt{1}: in un workflow di merge request, questo impedisce l'integrazione del codice finché le garanzie di sicurezza non sono soddisfatte. Il report prodotto rimane disponibile come artefatto scaricabile per l'analisi post-hoc; il gate decisionale è il codice di uscita.

La granularità attuale è binaria rispetto all'esito: o l'assessment è pulito o non lo è. Una traiettoria di sviluppo naturale è rendere il predicato di blocco configurabile per livello di priorità, così che una violazione di priorità P0 su un endpoint di autenticazione blocchi la pipeline mentre una violazione P3 su una configurazione subottimale produca un avviso senza interrompere il rilascio. Questa direzione è discussa nella Sezione~\ref{subsec:roadmap-operativa} come estensione operativa.

\subsection{Quality Gate a Due Velocità: Fail-Fast e Assessment Completo (D4.P8)}\label{subsec:fail-fast}

L'assessment completo descritto nel Capitolo~\ref{chap:experiments} ha richiesto circa cinque minuti su un target con 18 test attivi. In un gate su una merge request, dove l'obiettivo è verificare rapidamente che le modifiche non abbiano introdotto regressioni di sicurezza perimetrale, attendere il completamento dell'intero assessment è un costo sproporzionato rispetto alla decisione da prendere: se un test P0 rileva una violazione critica, il blocco della pipeline è già giustificato prima che i test successivi vengano eseguiti.

Il parametro \texttt{execution.fail\_fast} nel \texttt{config.yaml} introduce questa seconda dimensione di controllo, distinta dagli exit code ma complementare ad essi: mentre la Sezione~\ref{subsec:exit-codes} descrive cosa il tool comunica al termine dell'esecuzione, \texttt{fail\_fast} descrive se e quando il tool interrompe l'esecuzione prima del completamento naturale. Con il valore di default \texttt{false}, l'engine esegue tutti i test attivi e produce un report completo. Con \texttt{fail\_fast:~true}, l'engine interrompe la Phase~5 alla prima violazione confermata su un test di priorità P0: i test rimanenti non vengono eseguiti, il \texttt{ResultSet} contiene solo i risultati prodotti fino a quel punto, e il processo termina con exit~\texttt{1}. La condizione di interruzione include anche \texttt{ERROR} sui test P0, perché una verifica perimetrale incompleta per malfunzionamento del tool è operativamente equivalente a una violazione confermata: in entrambi i casi il controllo non garantisce ciò che deve garantire.

I due profili si adattano a scenari d'uso distinti. La modalità fail-fast è appropriata per i gate sulle merge request, dove la velocità di risposta ha valore: all'interruzione alla prima violazione P0, il wall-clock si riduce da cinque minuti a pochi secondi nel caso peggiore. La modalità completa è appropriata per assessment periodici schedulati su rami stabili, dove ogni dominio deve produrre evidenza per informare le decisioni di remediation. La combinazione con \texttt{execution.min\_priority} nel \texttt{config.yaml} aggiunge un'ulteriore dimensione: limitare l'esecuzione ai soli test P0 in modalità fail-fast produce un gate veloce sui controlli perimetrali critici; estendere a tutti i livelli di priorità in modalità completa produce l'assessment di profondità. Il medesimo tool, lo stesso file di configurazione di base, due comportamenti configurabili senza modifiche al codice.

\subsection{Variabilità di Privilegio nei Deployment Industriali}\label{subsec:privilegio-deployment}

Il tool adatta la propria copertura alle precondizioni disponibili senza che l'operatore scelga esplicitamente quale "modalità" attivare: la variazione è nell'insieme di credenziali fornite nel \texttt{config.yaml} e nel campo \texttt{gateway\_adapter}, e la degradazione verso livelli di copertura inferiori avviene in modo automatico e documentato, come stabilito nella Sezione~\ref{subsec:box-gradient}. Nei deployment industriali, questa proprietà si traduce in tre scenari concreti che riflettono le diverse condizioni di accesso delle organizzazioni che adottano il tool.

In contesti di integrazione continua su rami di sviluppo, i token di autenticazione per i ruoli configurati sono disponibili come segreti di pipeline, e la Admin \gls{API} del gateway è accessibile dall'interno della rete di staging. Tutti e tre i livelli del box gradient sono raggiungibili, e l'assessment copre l'intera matrice dei domini \texttt{BLACK\_BOX}, \texttt{GREY\_BOX} e \texttt{WHITE\_BOX}.

In contesti di verifica su ambienti di pre-produzione condivisi tra più team, le credenziali con privilegi amministrativi potrebbero non essere distribuite a tutti gli esecutori della pipeline per ragioni di governance interna. Il tool opera sui livelli \texttt{BLACK\_BOX} e \texttt{GREY\_BOX}: i test \texttt{WHITE\_BOX} producono \texttt{SKIP} con motivazione esplicita, e la riduzione di copertura è documentata nel report.

In contesti di red team o penetration test esterno, solo l'accesso di rete al target è garantito. Il tool opera in modalità \texttt{BLACK\_BOX}: i test \texttt{GREY\_BOX} e \texttt{WHITE\_BOX} producono \texttt{SKIP}, e l'assessment si concentra sui controlli perimetrali verificabili da un attaccante senza credenziali. In questo scenario, la copertura ridotta non è una limitazione da correggere: è il dato che misura la postura perimetrale del sistema così come la vedrebbe un attaccante esterno.


































% ─────────────────────────────────────────────────────────────────────────────
\section{Sviluppi Futuri}%
\label{sec:sviluppi-futuri}
% ─────────────────────────────────────────────────────────────────────────────

Le traiettorie di sviluppo si articolano su due scale temporali con natura diversa. Le estensioni operative sono conseguenze dirette dell'architettura corrente: connector progettati e non ancora implementati, domini predisposti nella tassonomia e non ancora coperti, componenti di pipeline che seguono direttamente dai meccanismi già validati. Non richiedono scelte di design nuove. Le traiettorie di ricerca pongono invece problemi aperti che richiedono analisi formale o validazione sperimentale indipendente prima dell'implementazione, e in alcuni casi costituiscono oggetti di studio autonomi. La distinzione non è di importanza ma di natura: confonderle produrrebbe impegni di consegna su lavoro che richiede ricerca.


\subsection{Estensioni Operative}\label{subsec:roadmap-operativa}

Le estensioni elencate in questa sezione seguono direttamente dall'architettura corrente: i punti di estensione sono già predisposti, i contratti tra i componenti sono già definiti. La mancata implementazione nella fase attuale riflette scelte di priorità, non vincoli architetturali. Le estensioni si raggruppano in tre aree.

\paragraph{Copertura dei test e connectors.}
Il Dominio~5 (Observability) è l'unico non coperto da nessun test nella fase attuale, come dichiarato nella Sezione~\ref{subsec:domini-tassonomia}. I test~5.1 (audit logging) e~5.2 (alert real-time su eventi anomali) sono già specificati nella metodologia; la loro implementazione richiede l'estensione del testbed con un log aggregator e un sistema di alerting nel Docker Compose, dipendenze di infrastruttura che si mantengono architetturalmente isolate dal core del tool.

Tra i connector Categoria~A non ancora implementati, il \texttt{JwtToolConnector} sblocca la copertura del ciclo di vita delle credenziali nel Dominio~1: i test~1.2 (validazione strutturale del \gls{JWT}) e~1.3 (expiry check), pianificati come \texttt{depends\_on~=~["1.1"]}, sono schedulati in una fase successiva del \gls{DAG} e non possono avviarsi finché test~1.1 non ha prodotto un token valido nel \texttt{TestContext}. La struttura del \gls{DAG} è già predisposta per accogliere questa Phase~C senza interventi sull'ordinamento topologico. Analogamente, il Dominio~2 (Authorization) nella fase attuale copre solo il test~2.1; i test~2.2--2.5, pianificati con il connector OFFAT per la generazione dinamica di probe \gls{BOLA}/\gls{IDOR} derivati dalla specifica OpenAPI, completerebbero la copertura del controllo accessi a livello di oggetto. Estendere il Dominio~3 con \texttt{Schemathesis} per l'injection testing\footnote{L'uso di Schemathesis come connector esterno, che genera e delega l'esecuzione dei probe, è distinto dall'uso come strumento di testing autonomo discusso nella Sezione~\ref{subsec:contract-driven}; qui Schemathesis opera come fornitore di dati grezzi, la cui valutazione rimane responsabilità del test.}, il Dominio~6 con \texttt{ext.6.3.socket} per HTTP request smuggling, e il Dominio~7 con connector Vegeta e Interactsh per race condition e \gls{SSRF} blind completano il quadro dei connector Cat~A pianificati.

Due di questi meritano una nota tecnica specifica. Il test~7.3 (\gls{TOCTOU}) richiede richieste simultanee con sincronizzazione sub-millisecondo: il \gls{GIL} di Python rende questo requisito inaffidabile, e il \texttt{VegetaConnector} delega la generazione del carico a Vegeta, scritto in Go con goroutine native, ottenendo la finestra di concorrenza necessaria. Il connector è riutilizzabile dal test~4.1 per il load testing del rate limiting, in coerenza con il principio di separazione dati/valutazione discusso nella Sezione~\ref{subsec:connector-hierarchy}. L'SSRF blind (test~\texttt{ext.7.4.interactsh}) richiede invece un oracle fuori banda: \texttt{InteractshConnector} registra un \gls{URL} controllato su un server interactsh remoto, lo inietta nel probe e verifica il callback ricevuto. L'evidenza non è nella risposta \gls{HTTP} al probe, ma nella richiesta in uscita dal server verso l'indirizzo controllato.

\paragraph{Integrazione della pipeline.}
Due estensioni riguardano il canale di comunicazione con la pipeline \gls{CI/CD}, già descritto nella Sezione~\ref{subsec:exit-codes}. La prima è la granularità degli exit code per livello di priorità, che permetterebbe di distinguere violazioni P0 (bloccanti in qualsiasi scenario) da violazioni P2/P3 (segnalabili come avvertimenti senza bloccare il rilascio). La seconda è l'implementazione degli adapter per i cloud provider principali (Amazon \gls{API} Gateway, Azure \gls{APIM}), che consentirebbe la copertura \texttt{WHITE\_BOX} nei deployment cloud managed discussi nella Sezione~\ref{subsec:limite-admin-api}.

\paragraph{Ottimizzazioni del motore.}
La struttura a batch del \gls{DAG} è già parallelizzabile per costruzione, come discusso nella Sezione~\ref{subsec:dag-impl}: i test all'interno di uno stesso batch non hanno dipendenze reciproche. Introdurre un \texttt{ThreadPoolExecutor} nella Phase~5 ridurrebbe il wall-clock proporzionalmente al numero di test eseguibili in parallelo, con un effetto stimabile dalla distribuzione osservata nella Sezione~\ref{sec:footprint-dag}. La ricerca necessaria riguarda i problemi di thread-safety su \texttt{EvidenceStore} e \texttt{TestContext}, discussi nella Sezione~\ref{subsec:ricerca-aperta}.


\subsection{Traiettorie di Ricerca Aperta}\label{subsec:ricerca-aperta}

Le direzioni descritte in questa sezione non sono estensioni implementative: sono problemi che l'architettura attuale abilita o porta alla luce, ma che richiedono analisi formale, validazione sperimentale indipendente o scelte di design nuove prima di poter essere affrontati concretamente.

\paragraph{Parallelizzazione intra-batch e thread-safety.}
L'estensione della parallelizzazione al livello intra-batch, anticipata nella Sezione~\ref{subsec:roadmap-operativa}, richiede la risoluzione di due problemi di thread-safety che non hanno soluzioni banali. Il primo riguarda l'\texttt{EvidenceStore}: il metodo \texttt{begin\_test()} apre un file in append e lo mantiene aperto per la durata del test; con più test concorrenti che chiamano \texttt{log\_transaction()}, le scritture devono essere serializzate senza introdurre un lock globale che annulli il vantaggio della parallelizzazione. Il secondo riguarda il \texttt{TestContext}: il canale dei token è condiviso tra tutti i test, e chiamate concorrenti a \texttt{acquire\_tokens()} potrebbero produrre login multipli per lo stesso ruolo, violando la proprietà di idempotenza discussa nella Sezione~\ref{subsec:auth-layer}. La strada percorribile, per entrambi i problemi, è la separazione dei contesti per batch anziché per run, con un \texttt{TestContext} locale a ogni batch e un meccanismo di merge controllato alla fine di ciascuno. L'analisi formale della correttezza di questo schema è il prerequisito all'implementazione.

\paragraph{APIGuard come piattaforma modulare.}
L'architettura attuale è costruita attorno al paradigma \gls{REST}/OpenAPI: i test nativi derivano la superficie di attacco dalla specifica \gls{OAS}, i connector integrano tool pensati per quel protocollo, e il modello di evidenza riflette la semantica di richiesta/risposta \gls{HTTP}. Esistono tuttavia classi di \gls{API} che questo modello non raggiunge: endpoint GraphQL, servizi gRPC, sistemi event-driven basati su AsyncAPI. La traiettoria di ricerca è trattare APIGuard Assurance non come un tool monoprotocollo ma come una piattaforma di assurance estendibile, dove il core (engine, \gls{DAG} scheduler, evidence store, report) rimane invariato e i moduli di protocollo vengono aggiunti come plugin indipendenti. Un modulo \texttt{apiguard-graphql} contribuirebbe i propri test nativi e connector per GraphQL rispettando i contratti di \texttt{BaseTest} e \texttt{BaseConnector}, senza modificare il core. Le sfide di ricerca riguardano la specifica formale del contratto di compatibilità tra moduli e core al variare delle versioni, la gestione della sicurezza dei moduli contribuiti da terze parti, e la definizione di un modello di evidenza comune abbastanza astratto da rappresentare semantiche di protocollo diverse.

\paragraph{Drift detection tra specifica e comportamento osservato.}
Il paradosso del Ground Truth discusso nella Sezione~\ref{subsec:limite-ground-truth}, per cui la specifica può essere obsoleta o incompleta senza che il tool possa rilevarlo, apre una traiettoria di ricerca che non richiede traffico passivo né inferenza statistica. Durante l'esecuzione, il tool raccoglie già osservazioni sistematiche sul comportamento del target: status code restituiti, header presenti o assenti, schemi di risposta effettivi. Confrontare formalmente queste osservazioni con le asserzioni della specifica OpenAPI potrebbe rivelare endpoint non documentati che rispondono, comportamenti difformi dalla specifica dichiarata, e silenziosamente produrre un indicatore di confidenza sulla completezza dell'oracolo usato. Questa direzione, diversa dall'inferenza di una specifica sintetica dal traffico, usa i dati che il tool già produce come input per una valutazione della qualità dell'oracolo stesso. Le sfide aperte riguardano la definizione formale di "difformità significativa", la distinzione tra deviazioni intenzionali e involontarie, e l'integrazione dell'indicatore nel report senza trasformarlo in un ulteriore vettore di falsi positivi.